\SourcePage{373}
\SGAchapter{XXI}{完全交的上同调的水平}{Le niveau de la cohomologie des intersections complètes}
\begin{center}
N. Katz\\
附 P. Deligne 撰写的附录（第 5 节）
\end{center}
% 原讲目录所列各节由译本目录重建；原目录页为PDF第373页。

\SourcePage{374}
\SGAsection{0}{引言}
设 $X$ 是\Term{dimension}{维数} $n\geq 3$ 的\Term{projective}{射影}\Term{smooth}{光滑}\Term{variety}{簇}，$Y(d)$ 是它与一个 $d$ 次\Term{hypersurface}{超曲面}的\Term{section}{截面}，$E^{n-1}(Y(d),\mathbf Q_\ell)$ 是 $Y(d)$ 的“\Term{vanishing-cohomology}{消失上同调}”（第 XVIII 讲 5.4.4）。本讲的目的是证明：当 $d$ 充分大且 $Y(d)$ 充分\Term{generic}{一般}时，并非整个 $E^{n-1}(Y(d),\mathbf Q_\ell)$ 都能通过\Term{algebraic-correspondence}{代数对应}“来自”\Term{proper}{固有}\Term{smooth}{光滑}\Term{algebraic-variety}{代数簇}的 $H^i$（$i<n-1$）。在\Term{characteristic-zero}{特征零}时，这是复\Term{cohomology}{上同调}的 Hodge \Term{bigrading}{双分次}以及\Term{coherent-cohomology}{凝聚上同调}满足下式的一个直接推论：
\[
 \dim H^{n-1}(Y(d),\mathcal O_{Y(d)})\longrightarrow\infty
 \qquad\text{当 }d\longrightarrow\infty.
\]
在\Term{characteristic}{特征} $p$ 时，所用的方法仍然是把\Term{coherent-cohomology}{凝聚上同调} $H^{n-1}(Y(d),\mathcal O_{Y(d)})$（一个\Term{characteristic}{特征} $p$ 的\Term{vector-space}{向量空间}）与\Term{etale-cohomology}{平展上同调} $H^{n-1}_{\text{ét}}(Y(d),\mathbf Q_\ell)$ 联系起来。为此，我们在一个\Term{finite-field}{有限}\Term{base-field}{基域}上工作；
\SourcePage{375}
在这种情形下，二者由簇 $Y(d)$ 的 \Term{zeta-function}{$\zeta$ 函数}联系起来：这一函数可用 $\ell$ 进\Term{cohomology}{上同调}表达，而它的“模 $p$ \Term{reduced}{约化}”可用\Term{coherent-cohomology}{凝聚上同调}表达（参见第 XXII 讲 3.9）。由于在\Term{complete-intersection}{完全交}的情形下，我们对这两种\Term{cohomology}{上同调}（\Term{etale-cohomology}{平展上同调}与\Term{coherent-cohomology}{凝聚上同调}）都有详细了解，这个方法可以顺利使用，并由关于\Term{hasse-witt-operation}{Hasse--Witt 运算}（1.0）的一个结果（1.4）\Term{pushout}{推出}主要结果 4.2。

A. Grothendieck 给出了一些依赖于\Term{crystalline-cohomology}{晶体上同调}理论的启发性论证，使人相信：对任意\Term{projective}{射影}\Term{smooth}{光滑} $X$，\Term{generic}{一般}命题也可以用本质上相同的方法得到。

即使在\Term{complete-intersection}{完全交}这一特殊情形下，证明也实质性地用到了（2.2）Frobenius \Term{eigenvalue}{特征值}的\Term{integrality}{整性}\Term{theorem}{定理} 2.1；该\Term{theorem}{定理}由 P. Deligne 最近证明，其证明见附录 I。

\SGAsection{1}{完全交的 Hasse--Witt 矩阵}
\begin{SGAresult}{定义}{1.0}
设 $k$ 是\Term{characteristic-exponent}{特征指数}为 $p$ 的域。一个 $k$ 上的\Term{proper}{固有}\Term{scheme}{概形} $X$ 称为\Term{special-scheme}{特殊的}（spécial），如果\Term{additive}{加性}的 \Term{p-linear}{$p$ 线性}运算
\[
 \mathfrak F_p:\mathcal O_X\longrightarrow\mathcal O_X,
 \qquad g\longmapsto g^p
 \tag{1.0.1}
\]
所诱导的运算
\[
 \mathfrak F_p:H^{\dim X}(X,\mathcal O_X)
 \longrightarrow H^{\dim X}(X,\mathcal O_X)
 \tag{1.0.2}
\]
\SourcePage{376}
是\Term{nilpotent}{幂零}的。\Term{equivalence}{等价}地，若令 $M=\dim_k H^{\dim X}(X,\mathcal O_X)$，则要求\Term{p-linear}{$p^M$ 线性}运算
\[
 \mathfrak F_p^M=\mathfrak F_{p^M}:\mathcal O_X\longrightarrow\mathcal O_X,
 \qquad g\longmapsto g^{p^M}
 \tag{1.0.3}
\]
在 $H^{\dim X}(X,\mathcal O_X)$ 上为零。
\end{SGAresult}
回顾文献 [2]：对每个 $i\in\mathbf Z$，由 (1.0.1) 诱导的\Term{homomorphism}{同态}
\[
 H^i(X,\mathcal O_X)\xrightarrow{\mathfrak F_p}H^i(X,\mathcal O_X)
\]
经典地称为 $H^i(X,\mathcal O_X)$ 中的\Term{hasse-witt-homomorphism}{Hasse--Witt 同态}。注意，当 $p=1$ 时，说 $X$ \Term{non-special-scheme}{非特殊}，仅仅意味着 $H^{\dim X}(X,\mathcal O_X)\ne0$。

\begin{SGAresult}{\Term{theorem}{定理}}{1.1}
固定一个\Term{characteristic-exponent}{特征指数} $p$、一个\Term{dimension}{维数} $n\geq1$、一个整数 $r\geq1$ 及一个\Term{multidegree}{多重次数} $(d_1,\ldots,d_r)$。若
\[
 \sum_{i=1}^r d_i\geq(n+r)+1,
 \tag{1.1.0}
\]
则在\Term{characteristic}{特征} $p$ 时，$\mathbf P^{n+r}$ 中\Term{multidegree}{多重次数}为 $(d_1,\ldots,d_r)$ 的所有（$n$ 维）\Term{complete-intersection}{完全交}中，一个“充分\Term{generic}{一般}”的成员是\Term{non-special-scheme}{非特殊}的（1.0）。
\end{SGAresult}
\begin{proof}
记
\[
 f:X_U\longrightarrow U
 \tag{1.1.1}
\]
为\Term{characteristic}{特征} $p$ 时 $\mathbf P^{n+r}$ 中\Term{multidegree}{多重次数}为 $(d_1,\ldots,d_r)$ 的\Term{complete-intersection}{完全交}的“\Term{universal-family}{泛族}”
\SourcePage{377}
（带有其定义方程）；于是 $U$ 是 $\mathbf F_p$ 上 $r$ 个\Term{projective-space}{射影空间}之积中的一个\Term{open-subscheme}{开子概形}，其\Term{dimension}{维数}相当可观。

我们知道，$\mathcal O_U$ 模
\[
 R^n f_*\mathcal O_{X_U}
 \tag{1.1.2}
\]
是\Term{locally-free}{局部自由}的，且其形成与任意\Term{base-change}{基变换}交换（第 XI 讲 1.5(i)）；由假设 (1.1.0)，其\Term{rank}{秩} $M$ 非零（第 XI 讲 2.5）。当 $p=1$ 时，这就对每个满足 (1.1.0) 的\Term{complete-intersection}{完全交}证明了\Term{theorem}{定理}。

设 $p\geq2$。迭代的\Term{hasse-witt-operation}{Hasse--Witt 运算}
\[
 \mathfrak F_p^M=\mathfrak F_{p^M}:
 R^n f_*\mathcal O_{X_U}\longrightarrow R^n f_*\mathcal O_{X_U}
 \tag{1.1.3}
\]
在 $\mathcal O_U$ 上是\Term{p-linear}{$p^M$ 线性}的，又因 $R^n f_*\mathcal O_{X_U}$ \Term{locally-free}{局部自由}，所以运算 (1.1.3) 在 $U$ 上\Term{local-scheme}{局部}可由一个取值于 $\mathcal O_U$ 的 $M\times M$ 矩阵 $H$ 表示。因此
\[
 \{u\in U\mid X_u\text{ 非特殊}\}\text{ 是 }U\text{ 中的开集}
 \tag{1.1.4}
\]
（因为它就是 $H(u)\ne0$ 的那些 $u$ 的\Term{set}{集合}）。又因 $U$ \Term{irreducible}{不可约}，只需找到一个使 $X_u$ \Term{non-special-scheme}{非特殊}的 $u$；我们将在\Term{finite-field}{有限域}上做到这一点。
\end{proof}

\begin{SGAresult}{命题}{1.1.5}
设 $X$ 是 $\mathbf F_q$ 上的一个\Term{complete-intersection}{完全交}。若 $X$ 没有 $\mathbf F_q$ \Term{rational-point}{有理点}，则 $X$ \Term{non-special-scheme}{非特殊}。
\end{SGAresult}
\SourcePage{378}
\begin{proof}
显然，只需证明对\Term{complete-intersection}{完全交} $X$ 有如下蕴含：
\[
 X\text{ 特殊}\quad\Longrightarrow\quad
 \operatorname{card}X(\mathbf F_q)\equiv1\pmod p.
 \tag{1.1.5.1}
\]
由\Term{trace-formula}{迹公式}（在下一讲，第 XXII 讲 3.2 中证明），有
\[
 \operatorname{card}X(\mathbf F_q)\equiv
 \sum_{i\geq0}(-1)^i\operatorname{tr}
 \bigl(\mathfrak F_q\mid H^i(X,\mathcal O_X)\bigr)\pmod p.
\]
考虑到 $i\ne0,n=\dim X$ 时 $H^i(X,\mathcal O_X)=0$，且 $H^0(X,\mathcal O_X)\simeq\mathbf F_q$（第 XI 讲 1.5(iii)），可将此式写成
\[
 \operatorname{card}X(\mathbf F_q)\equiv
 1+(-1)^n\operatorname{tr}\bigl(\mathfrak F_q\mid H^n(X,\mathcal O_X)\bigr)\pmod p.
\]
若 $X$ \Term{special-scheme}{特殊}，则依定义，$\mathfrak F_q$ 在 $H^n(X,\mathcal O_X)$ 上的\Term{action}{作用}\Term{nilpotent}{幂零}，所以其迹为零。这就得到蕴含 (1.1.5.1)，从而得到命题。
\end{proof}

\begin{SGAresult}{命题}{1.2}
无论 $p$ 的幂 $q$ 如何选取，只要
\[
 \sum_{i=1}^r d_i\geq n+r+1,
\]
就存在 $\mathbf P^{n+r}_{\mathbf F_q}$ 中一个\Term{multidegree}{多重次数}为 $(d_1,\ldots,d_r)$、没有 $\mathbf F_q$ \Term{rational-point}{有理点}的\Term{complete-intersection}{完全交}。
\end{SGAresult}
\begin{proof}
先看 $r=1$ 的情形。对每个 $d\geq n+2$，需找一个 $d$ 次\Term{homogeneous-form}{齐次式}
\SourcePage{379}
\[
 H(X_1,\ldots,X_{n+2})\in\mathbf F_q[X_1,\ldots,X_{n+2}]
 \tag{1.2.1}
\]
使得
\[
 \begin{gathered}
 (\alpha_1,\ldots,\alpha_{n+2})\in\mathbf F_q^{n+2},\quad
 H(\alpha_1,\ldots,\alpha_{n+2})=0\\
 \Longrightarrow\quad\alpha_1=\cdots=\alpha_{n+2}=0.
 \end{gathered}
 \tag{1.2.2}
\]
为此，令 $E=\mathbf F_{q^d}$，并取 $E$ 中在 $\mathbf F_q$ 上\Term{linearly-independent}{线性无关}的元素 $\beta_1,\ldots,\beta_{n+2}$。取
\[
 H(X_1,\ldots,X_{n+2})=
 \operatorname N_{E/\mathbf F_q}\left(\sum_{i=1}^{n+2}\beta_iX_i\right).
 \tag{1.2.3}
\]
在\Term{generic}{一般}情形下，由于
\[
 \sum_{i=1}^r d_i\geq n+r+1,
 \tag{1.2.4}
\]
可以找到 $r$ 个整数
\[
 1\leq a_1<a_2<\cdots<a_r=n+r+1
 \tag{1.2.5}
\]
满足
\[
 \left\{\begin{aligned}
 d_1&\geq a_0,\\
 d_2&\geq a_2-a_1,\\
 &\ \vdots\\
 d_r&\geq a_r-a_{r-1}.
 \end{aligned}\right.
 \tag{1.2.6}
\]
\SourcePage{380}
对 $i=1,\ldots,r$，取一个关于 $a_i-a_{i-1}$ 个变量的 $d_i$ 次\Term{homogeneous-form}{齐次式}
\[
 H_i(X_{a_{i-1}+1},\ldots,X_{a_i})
 \in\mathbf F_q[X_{a_{i-1}+1},\ldots,X_{a_i}],
\]
使之具有性质 (1.2.2)。所求的\Term{complete-intersection}{完全交}就是由\Term{homogeneous-form}{齐次式}
\[
 H_1(X_1,\ldots,X_{a_1}),\quad\ldots,\quad
 H_r(X_{a_{r-1}},\ldots,X_{a_r})
\]
定义的那个。
\end{proof}

\begin{SGAresult}{\Term{theorem}{定理}}{1.3}
设 $X$ 是\Term{characteristic-exponent}{特征指数}为 $p$ 的域 $k$ 上一个\Term{dimension}{维数} $n\geq1$ 的\Term{irreducible}{不可约}\Term{complete-intersection}{完全交}。对每个整数 $d\geq1+n$，$X$ 与一个充分\Term{generic}{一般}的 $d$ 次\Term{hypersurface}{超曲面}的\Term{section}{截面}是\Term{non-special-scheme}{非特殊}的（1.0）。
\end{SGAresult}
\begin{proof}
$X$ 与 $d$ 次（\Term{hypersurface}{超曲面}）的\Term{section}{截面} $X.H_u$ 由 $H^0(X,\mathcal O_X(d))$ 中的直线参数化，即由
\[
 \mathbf P\bigl(H^0(X,\mathcal O_X(d))\bigr)=U
 \tag{1.3.1}
\]
参数化。在 $U$ 上，我们有“$X$ 与一个 $d$ 次\Term{hypersurface}{超曲面}的\Term{universal-section}{泛截面}”
\[
 f:X.H_U\longrightarrow U.
 \tag{1.3.2}
\]
设 $U'$ 是由满足 $\dim X_u=n-1$ 的 $u\in U$ 组成的\Term{open-subset}{开集}。由第 XI 讲 1.5(i) 和 2.5，我们知道，$\mathcal O_U$ 模
\[
 R^{n-1}f_*\mathcal O_{X.H_U}\bigm|_{U'}
 \tag{1.3.3}
\]
\SourcePage{381}
是一个\Term{locally-free}{局部自由}\Term{sheaf}{层}，其形成与任意\Term{base-change}{基变换}交换，且其\Term{rank}{秩} $M\geq1$。这就在 $p=1$ 时，对每个与\Term{degree}{次数} $\geq1+n$ 的\Term{hypersurface}{超曲面}所得的\Term{section}{截面}证明了\Term{theorem}{定理}。因此设 $p>0$；运算 $\mathfrak F_p^n$ 仍然可以\Term{local-scheme}{局部}表示为取值于 $\mathcal O_U$ 的一个 $M\times M$ 矩阵，因而
\[
 \{u\in U'\mid X.H_u\text{ 非特殊}\}\text{ 是 }U'\text{ 中的开集}.
 \tag{1.3.4}
\]
所以只需证明\Term{set}{集合} (1.3.4) \Term{nonempty}{非空}。

先处理
\[
 k=\mathbf F_q
 \tag{1.3.5}
\]
的情形。将 $\mathbf F_q$ 换成一个\Term{finite-extension}{有限扩张}后，可设
\[
 \operatorname{card}X(\mathbf F_q)>0
 \tag{1.3.6}
\]
以及
\[
 \begin{gathered}
 \text{环境空间 }\mathbf P^{n+r}_k\text{ 中存在齐次坐标 }X_1,\ldots,X_{n+r+1},\\
 \text{使 }X\text{ 不与线性子簇 }X_1=\cdots=X_{n+1}=0\text{ 相交}.
 \end{gathered}
 \tag{1.3.7}
\]
现在取一个 $d$ 次\Term{homogeneous-form}{齐次式}
\[
 H(X_1,\ldots,X_{n+1})\in\mathbf F_q[X_1,\ldots,X_{n+1}]
 \tag{1.3.8}
\]
使之具有性质 (1.2.2)（因 $d\geq1+n$，这是可能的）。取 $H$ 为 $\mathbf P^{n+r}$ 中方程 $H(X_1,\ldots,X_{n+1})=0$ 定义的\Term{hypersurface}{超曲面}。
\SourcePage{382}
按构造，交 $X.H$ 没有 $\mathbf F_q$ \Term{rational-point}{有理点}。结合 (1.3.6)，这同时说明 $X\not\subset H$，故 $X.H$ 的\Term{dimension}{维数}为 $n-1$（因为 $X$ \Term{irreducible}{不可约}），即 $H\in U'$；也说明 $X.H$ \Term{non-special-scheme}{非特殊}（1.5）。

再转向\Term{generic}{一般}情形。考察定义 $X$ 的方程，可以看出可设 $k$ 在 $\mathbf F_p$ 上\Term{finite-type}{有限型}，并存在 $k$ 的\Term{subring}{子环} $B$，它作为环在 $\mathbf F_p$ 上\Term{finite-type}{有限型}，\Term{fraction-field}{分式域}为 $k$，满足
\[
 \begin{gathered}
 X\subset\mathbf P^{n+r}_k\text{ 是一个相对完全交 }
 \widetilde X\subset\mathbf P^{n+r}_B\text{ 的泛纤维，}\\
 \text{且后者逐纤维与 }X\text{ 有相同的多重次数，并且几何不可约}.
 \end{gathered}
 \tag{1.3.9}
\]
记\Term{projection}{投影}为
\[
 \rho:\widetilde X\longrightarrow S=\operatorname{Spec}B,
 \tag{1.3.10}
\]
记
\[
 E\text{ 为局部自由的 }\mathcal O_S\text{ 模 }\rho_*\mathcal O_{\widetilde X}(d),
 \tag{1.3.11}
\]
并记
\[
 \pi:T=\mathbf P^\vee(E)\longrightarrow S
 \tag{1.3.12}
\]
为与 $E$ 中直线相应的\Term{bundle}{丛}。注意，(1.3.12) 的\Term{generic-fibre}{泛纤维}正是 $U=\mathbf P(H^0(X,\mathcal O_X(d)))$。在 $T$ 上，有 $\widetilde X\subset\mathbf P^{n+r}_B$ 与 $d$ 次\Term{hypersurface}{超曲面}的\Term{universal-section}{泛截面}，记为
\SourcePage{383}
\[
 \alpha:\widetilde X.H_T\longrightarrow T.
 \tag{1.3.13}
\]
这些构造组成如下各方块均为\Term{cartesian-square}{笛卡尔方块}的图：
\[
\begin{tikzcd}[column sep=large,row sep=large]
 \widetilde X.H_T\arrow[d,"\alpha"] &
 X.H_U=(\widetilde X.H_T)\times_U T\arrow[l]\arrow[d]\\
 T\arrow[d,"\pi"] & U=T\times_S\operatorname{Spec}k\arrow[l]\arrow[d]\\
 S & \operatorname{Spec}k\arrow[l]
\end{tikzcd}
\tag{1.3.14}
\]
而且，$\mathcal O_T$ 模
\[
 R^{n-1}\alpha_*\mathcal O_{\widetilde X.H_T}
 \tag{1.3.15}
\]
在由满足 $\dim\alpha^{-1}(t)=n-1$ 的 $t$ 所组成的\Term{open-subset}{开集} $T'\subset T$ 上仍是\Term{locally-free}{局部自由}的，且其形成在那里与任意\Term{base-change}{基变换}交换。由于\Term{scheme}{概形} $T$ 与 $U$ 有相同的\Term{generic-point}{泛点}，要证明 (1.3.4) \Term{nonempty}{非空}，只需验证 $T'$ 的\Term{open-subset}{开集}
\[
 \{t\in T'\mid\widetilde X_{\pi(t)}.H_t\text{ 非特殊}\}
 \tag{1.3.16}
\]
\Term{nonempty}{非空}。已经处理的情形 (1.3.5) 说明：对每个
\SourcePage{384}
\Term{closed-point}{闭点} $s\in S$，\Term{fibre}{纤维} $\pi^{-1}(s)$ 都与\Term{open-subset}{开集} (1.3.16) \Term{nonempty}{非空}相交。证毕。
\end{proof}

完全类似的论证可以证明：
\begin{SGAresult}{\Term{theorem}{定理}}{1.4}
设 $X\subset\mathbf P^N$ 是域 $k$ 上一个\Term{equidimensional}{等维}的、\Term{dimension}{维数} $n\geq1$ 的\Term{projective}{射影}\Term{cohen-macaulay}{Cohen--Macaulay}\Term{scheme}{概形}。则存在整数 $d_0$，使得当 $d\geq d_0$ 时，若以 $X.H\xrightarrow{i}X$ 表示与一个\Term{generic}{一般} $d$ 次\Term{hypersurface}{超曲面}所得\Term{section}{截面}的\Term{inclusion}{包含映射}，则
\[
 \mathfrak F_p\text{ 在 }
 H^{n-1}(X.H,\mathcal O_{X.H})/i^*H^{n-1}(X,\mathcal O_X)
 \text{ 上的作用}
 \tag{1.4.0}
\]
不是\Term{nilpotent}{幂零}的。
\end{SGAresult}
\begin{proof}
如上所述，可归结到 $k=\mathbf F_q$ 的情形，并归结为寻找一个使 (1.4.0) 成立的 $H$（\Term{degree}{次数}给定且充分大，并使 $\dim H^{n-1}(X.H,\mathcal O_{X.H})$ 与 $H$ 无关；参见 1.4.4）。将 $q$ 扩大后，可设 $X$ 的\Term{irreducible-component}{不可约分支} $X_i$ 均\Term{geometrically-irreducible}{几何不可约}，满足
\[
 \operatorname{Card}X_i(\mathbf F_q)\geq2,
 \tag{1.4.1}
\]
并且
\[
 \begin{gathered}
 \text{存在射影坐标 }X_0,\ldots,X_N,\text{ 使 }X\text{ 不与}\\
 \text{线性子簇 }X_0=\cdots=X_n=0\text{ 相交}.
 \end{gathered}
 \tag{1.4.2}
\]
由 Serre 的一个\Term{vanishing-theorem}{消失定理}（[1]，或 SGA 2 第 XII 讲 1.4），存在一个整数 $d_1$，使每个整数 $d\geq d_1$ 都满足
\SourcePage{385}
\[
 H^i(X,\mathcal O_X(-d))=0\qquad\text{若 }i\ne n=\dim X.
 \tag{1.4.3}
\]
以 $H$ 表示一个\Term{degree}{次数} $d\geq d_1$ 的\Term{hypersurface}{超曲面}，由 (1.4.3) 有
\[
 \begin{aligned}
 H^i(X,\mathcal O_X)&\xrightarrow{\sim}H^i(X.H,\mathcal O_{X.H})
 &&\text{若 }i\leq n-2,\\
 H^{n-1}(X,\mathcal O_X)&\lhook\joinrel\longrightarrow
 H^{n-1}(X.H,\mathcal O_{X.H})&&\text{为单射}.
 \end{aligned}
 \tag{1.4.4}
\]
我们将取 $d_0=\max(1+N,d_1)$ 来证明\Term{theorem}{定理}。对每个整数 $d\geq d_0$，需找到一个定义在 $\mathbf F_q$ 上的 $d$ 次\Term{hypersurface}{超曲面} $H$，它不包含任何 $X_i$，并且
\[
 \operatorname{trace}\left(\mathfrak F_q\mid
 H^{n-1}(X.H,\mathcal O_{X.H})/H^{n-1}(X,\mathcal O_X)\right)\ne0.
 \tag{1.4.5}
\]
由\Term{trace-formula}{迹公式}（第 XX 讲 3.2.1），有
\[
\begin{aligned}
 \operatorname{card}X.H(\mathbf F_q)\bmod p
 &=\sum_{i=0}^{n-1}(-1)^i
 \operatorname{tr}\bigl(\mathfrak F_q\mid H^i(X.H,\mathcal O_{X.H})\bigr)\\
 &\overset{(1.4.4)}{=}
 \sum_{i=0}^{n-1}(-1)^i\operatorname{tr}
 \bigl(\mathfrak F_q\mid H^i(X,\mathcal O_X)\bigr)\\
 &\quad +(1)^{n-1}\operatorname{tr}\left(\mathfrak F_q\mid
 \frac{H^{n-1}(X.H,\mathcal O_{X.H})}{H^{n-1}(X,\mathcal O_X)}\right).
\end{aligned}
\tag{1.4.6}
\]
以 $a\in\mathbf Z/p\mathbf Z$ 表示量
\[
 \sum_{i=0}^{n-1}(-1)^i\operatorname{tr}
 \bigl(\mathfrak F_q\mid H^i(X,\mathcal O_X)\bigr).
\]
\SourcePage{386}
于是对每个 $d\geq d_0$，只需找到一个定义在 $\mathbf F_q$ 上、不包含 $X$ 的 $d$ 次 $H$，满足
\[
 \operatorname{Card}X.H(\mathbf F_q)\not\equiv a\pmod p.
 \tag{1.4.7}
\]
应分为两种情形。若 $a\ne0$，只需如 (1.3.8) 那样，取一个 $\mathbf F_q$ 上的 $H$，使 $(X.H)(\mathbf F_q)$ 为空。这样的 $H$ 满足 (1.4.7)，且不包含任何 $X_i$（因为依假设，$\operatorname{card}X_i(\mathbf F_q)\geq2$）。

若 $a=0$，只需找到一个 $H$，使 $\operatorname{Card}(X.H)(\mathbf F_q)=1$（这就满足了 (1.4.7)，且 $X_i\not\subset H$，因为 $\operatorname{Card}X_i(\mathbf F_q)\geq2$）。为此，选择\Term{projective-coordinate}{射影坐标} $X_0,\ldots,X_N$，使 $(1,0,\ldots,0)$ 是 $X$ 的一个点。取 $H$ 的方程为 $H(X_1,\ldots,X_N)=0$，其中 $H$ 是满足下述条件的\Term{homogeneous-form}{齐次式}（参见 (1.2.2)）：
\[
 \alpha_1,\ldots,\alpha_N\in\mathbf F_q,\quad
 H(\alpha_1,\ldots,\alpha_N)=0
 \quad\Longrightarrow\quad\alpha_1=\cdots=\alpha_N=0.
 \tag{1.4.8}
\]
于是 $H$ 在 $\mathbf P^N$ 中定义的\Term{hypersurface}{超曲面}只有 $(1,0,\ldots,0)$ 这一个 $\mathbf F_q$ 值点，因而 $X.H$ 也只有这一个 $\mathbf F_q$ 值点。证毕。
\end{proof}

\SourcePage{387}
\SGAsection{2}{定义域有限时的余水平}
\textbf{2.0.} 设 $X$ 是 $\mathbf F_q$ 上\Term{separated}{分离}的\Term{finite-type}{有限型}\Term{scheme}{概形}。对每个素数 $\ell\ne\operatorname{car}\mathbf F_q$，记
\[
 H^i_c(\overline X,\mathbf Q_\ell)
 \tag{2.0.1}
\]
为 $\overline X=X\otimes_{\mathbf F_q}\overline{\mathbf F}_q$ 的\Term{proper-support-cohomology}{固有支撑上同调}群。它们是\Term{finite}{有限}\Term{rank}{秩}的 $\mathbf Q_\ell$ \Term{vector-space}{向量空间}，$\mathbf F_q$ 的\Term{galois-group}{Galois 群} $\operatorname{Gal}(/\mathbf F_q)$ 在其上\Term{action}{作用}。记
\[
 \varphi_q\in\operatorname{Gal}(/\mathbf F_q)
 \quad\text{为 Frobenius 自同构}.
 \tag{2.0.2}
\]
我们知道（SGA 5 第 XIII 讲 3.2.1），$X$ 的\Term{zeta-function}{zeta 函数}（第 XXII 讲 3.0.1）由下式给出：
\[
 Z(t,X/\mathbf F_q)=\prod_{i\geq0}
 \det\bigl(1-t\varphi_q^{-1}\mid H^i_c(\overline X,\mathbf Q_\ell)\bigr)^{(-1)^{i+1}}.
 \tag{2.0.3}
\]
Deligne 最近将 Washnitzer 在 $X$ \Term{projective}{射影}且\Term{smooth}{光滑}时得到的结果推广，证明了如下\Term{theorem}{定理}；其证明见下文附录（5、5.3）：
\begin{SGAresult}{\Term{theorem}{定理}}{2.1}
$\varphi_q^{-1}$ 在 $H^i_c(\overline X,\mathbf Q_\ell)$ 中的\Term{eigenvalue}{特征值}均为\Term{algebraic-integer}{代数整数}。若 $i>\dim X$，则它们作为\Term{algebraic-integer}{代数整数}可被 $q^{i-\dim X}$ 整除。
\end{SGAresult}
\SourcePage{388}
特别地，直接得到：
\begin{SGAresult}{}{2.1 bis}
$\varphi_q^{-1}$ 在 $H^i_c(\overline X,\mathbf Q_\ell(-j))$ 中的\Term{eigenvalue}{特征值}是可被 $q^j$ 整除的\Term{algebraic-integer}{代数整数}。
\end{SGAresult}
\begin{SGAresult}{推论}{2.2}
设 $X$ 是 $\mathbf F_q$ 上\Term{proper}{固有}、\Term{smooth}{光滑}且\Term{geometrically-connected}{几何连通}的\Term{scheme}{概形}。则 $\varphi_q^{-1}$ 在 $N^jH^q(\overline X,\mathbf Q_\ell)$ 中的\Term{eigenvalue}{特征值}是可被 $q^j$ 整除的\Term{algebraic-integer}{代数整数}（记号见第 XX 讲 2.0.6）。
\end{SGAresult}
\begin{proof}
该命题在\Term{finite-extension}{有限扩张} $\mathbf F_{q^M}/\mathbf F_q$ 下\Term{invariant}{不变}（这相当于将所考虑的\Term{eigenvalue}{特征值}取 $M$ 次幂）。作这样的\Term{base-change}{基变换}后，依定义，有\Term{finite}{有限}多个 $\mathbf F_q$ 上的\Term{projective}{射影}、\Term{smooth}{光滑}且\Term{geometrically-connected}{几何连通}\Term{scheme}{概形} $T_i$，其\Term{dimension}{维数}为 $d_i$，并有\Term{algebraic-cycle}{代数圈}
\[
 Z_i\text{ 在 }T_i\times_{\mathbf F_q}X\text{ 上，余维为 }j+d_i,
\]
使\Term{galois-group}{Galois 群} $\operatorname{Gal}(/\mathbf F_q)$ 模的\Term{homomorphism}{同态}
\[
 H^{q-2j}(\overline T_i,\mathbf Q_\ell(-j))
 \xrightarrow{\ \bigoplus\operatorname{cl}(Z_i)\ }
 N^jH^q(\overline X,\mathbf Q_\ell)
\]
为\Term{surjective}{满射}。由 2.1 bis 即得结论。
\end{proof}

\SourcePage{389}
\SGAsection{3}{对完全交的应用}
\textbf{3.0.} 设 $X$ 是域 $k$ 上 $\mathbf P^N$ 中一个\Term{dimension}{维数} $n\geq1$ 的\Term{smooth}{光滑}\Term{complete-intersection}{完全交}。对 $\ell\ne\operatorname{car}k$，由“\Term{weak-lefschetz}{弱 Lefschetz 定理}”和\Term{duality}{对偶性}，我们知道（第 XI 讲 1.6）
\[
 H^i(\overline X,\mathbf Q_\ell)=
 \begin{cases}
 0&\text{若 }i\text{ 为奇数且 }i\ne n,\\
 \mathbf Q_\ell(-i/2)&\text{若 }i\text{ 为偶数、}i\leq2n\text{ 且 }i\ne n,\\
 H^n(\overline X,\mathbf Q_\ell)&\text{若 }i=n.
 \end{cases}
 \tag{3.0.1}
\]
当 $k=\mathbf F_q$ 时，由此得到\Term{zeta-function}{zeta 函数}的公式：
\[
 Z(t,X/\mathbf F_q)\prod_{i=0}^n(1-q^it)
 =\det\bigl(1-t\varphi_q^{-1}\mid
 \operatorname{Prim}^n(\overline X,\mathbf Q_\ell)\bigr)^{(-1)^{n+1}}.
 \tag{3.0.2}
\]
（这里用了\Term{primitive-decomposition}{本原分解}和 (3.0.1)。）
\begin{SGAresult}{推论}{3.0.3}
在 3.0 的假设下，若 $k=\mathbf F_q$，则有模 $p$ 的同余式
\[
 \det\bigl(1-t\varphi_q^{-1}\mid H^n(\overline X,\mathbf Q_\ell)\bigr)
 \equiv\det\bigl(1-t\mathfrak F_q\mid H^n(X,\mathcal O_X)\bigr)\pmod p.
\]
\end{SGAresult}
\begin{proof}
由同余公式（第 XXII 讲 3.2.1），并考虑到
\[
 H^i(X,\mathcal O_X)\simeq
 \begin{cases}
 \mathbf F_q&i=0,\\
 0&1\leq i\leq n-1,\\
 0&i\geq n+1,
 \end{cases}
 \tag{3.0.4}
\]
\SourcePage{390}
有
\[
 Z(t,X/\mathbf F_q)\equiv(1-t)^{-1}
 \det\bigl(1-t\mathfrak F_q\mid H^n(X,\mathcal O_X)\bigr)^{(-1)^{n+1}}
 \pmod p.
 \tag{3.0.5}
\]
因此，由 (3.0.5) 和 (3.0.2) 得到同余式
\[
 \begin{aligned}
 {}&\det\bigl(1-t\varphi_q^{-1}\mid
 \operatorname{Prim}^n(\overline X,\mathbf Q_\ell)\bigr)\\
 &\hspace{2em}\equiv\det\bigl(1-t\mathfrak F_q\mid H^n(X,\mathcal O_X)\bigr)\pmod p.
 \end{aligned}
 \tag{3.0.6}
\]
所需结论由此得到，因为
\[
 H^n(\overline X,\mathbf Q_\ell)\simeq
 \begin{cases}
 \operatorname{Prim}^n(\overline X,\mathbf Q_\ell)&\text{若 }n\text{ 为奇数},\\
 \operatorname{Prim}^n(\overline X,\mathbf Q_\ell)\oplus\mathbf Q_\ell(-n/2)
 &\text{若 }n\text{ 为偶数}.
 \end{cases}
 \tag{3.0.7}
\]
\end{proof}

\begin{SGAresult}{推论}{3.1}
设 $X$ 是 $\mathbf F_q$ 上一个\Term{dimension}{维数} $n\geq1$ 的\Term{smooth}{光滑}\Term{complete-intersection}{完全交}。若
\[
 H^n(\overline X,\mathbf Q_\ell)=N^1H^n(\overline X,\mathbf Q_\ell),
 \tag{3.1.0}
\]
则
\[
 \mathfrak F_p\text{ 在 }H^n(X,\mathcal O_X)\text{ 上幂零}.
 \tag{3.1.1}
\]
\end{SGAresult}
\begin{proof}
由 (2.2)，(3.1.0) 蕴含
\[
 \det\bigl(1-t\varphi_q^{-1}\mid H^n(\overline X,\mathbf Q_\ell)\bigr)
 \equiv1\pmod p\quad\text{（甚至模 }q\text{）}.
 \tag{3.1.2}
\]
结合 (3.0.3)，这蕴含
\[
 \det\bigl(1-t\mathfrak F_q\mid H^n(X,\mathcal O_X)\bigr)=1,
 \tag{3.1.3}
\]
\SourcePage{391}
于是 $\mathfrak F_q$ 在 $H^n(X,\mathcal O_X)$ 上\Term{nilpotent}{幂零}，故 $\mathfrak F_p$ 也\Term{nilpotent}{幂零}。证毕。
\end{proof}

\begin{SGAresult}{命题}{3.2}
设 $X$ 是域 $k$ 上一个\Term{dimension}{维数} $n\geq1$ 的\Term{smooth}{光滑}\Term{complete-intersection}{完全交}。若
\[
 H^n(\overline X,\mathbf Q_\ell)=N^1H^n(\overline X,\mathbf Q_\ell),
 \tag{3.2.1}
\]
则 $X$ \Term{special-scheme}{特殊}，即
\[
 \mathfrak F_p\text{ 在 }H^n(X,\mathcal O_X)\text{ 上幂零}.
 \tag{3.2.2}
\]
\end{SGAresult}
\begin{proof}
可设 $k$ 在 $\mathbf F_p$ 上\Term{finite-type}{有限型}（$p=0$ 的情形是显然的），并且 $k$ 中存在一个\Term{finite-type}{有限型} $\mathbf F_p$ \Term{subalgebra}{子代数} $B$，使 $X/k$ 是一个相对\Term{smooth}{光滑}\Term{complete-intersection}{完全交} $\widetilde X$ 的\Term{generic-fibre}{泛纤维}：
\[
\begin{tikzcd}[column sep=large,row sep=large]
 \widetilde X\arrow[d,"\pi"]&X=\widetilde X_\eta\arrow[l]\arrow[d]\\
 \operatorname{Spec}B=S&\eta=\operatorname{Spec}k\arrow[l]
\end{tikzcd}
\tag{3.2.3}
\]
我们知道，$B$ 模 $R^n\pi_*\mathcal O_{\widetilde X}$ \Term{locally-free}{局部自由}，且其形成与任意\Term{base-change}{基变换}交换。因此点集
\SourcePage{392}
\[
 \{s\in S\mid X_s\text{ 特殊}\}
 \tag{3.2.4}
\]
是 $S$ 的一个\Term{closed-subset}{闭集}。

由第 XX 讲 2.3，将 $S$ 缩小后，可设
\[
 \forall s\in S,\qquad
 H^n(X_{\overline s},\mathbf Q_\ell)=N^1H^n(X_{\overline s},\mathbf Q_\ell).
 \tag{3.2.5}
\]
由 (3.1)，对 $S$ 的每个\Term{closed-point}{闭点} $s$，$X_s$ 都\Term{special-scheme}{特殊}。从而 $S$ 的\Term{closed-subset}{闭集} (3.2.4) 就是整个 $S$；特别地，$X=X_\eta$ \Term{special-scheme}{特殊}。
\end{proof}

\SGAsection{4}{结论}
结合 1.1 和 3.2，得到：
\begin{SGAresult}{\Term{theorem}{定理}}{4.1}
固定一个\Term{characteristic}{特征} $p$、一个\Term{dimension}{维数} $n$ 及一个\Term{multidegree}{多重次数} $(d_1,\ldots,d_r)$。记
\[
 X_V\longrightarrow V
 \tag{4.1.0}
\]
为\Term{characteristic}{特征} $p$ 时\Term{dimension}{维数}为 $n$、\Term{multidegree}{多重次数}为 $(d_1,\ldots,d_r)$ 的\Term{universal-family}{泛}\Term{smooth}{光滑}\Term{complete-intersection}{完全交}。若
\[
 \sum_{i=1}^r d_i\geq n+r+1,
 \tag{4.1.1}
\]
则存在一个\Term{dense-open}{稠密开集} $U\subset V$，使得
\[
 \text{对 }u\in U,\qquad
 H^n(X_{\overline u},\mathbf Q_\ell)\ne N^1H^n(X_{\overline u},\mathbf Q_\ell).
 \tag{4.1.2}
\]
\end{SGAresult}

\SourcePage{393}
同样，结合 1.3 和 3.2，得到：
\begin{SGAresult}{\Term{theorem}{定理}}{4.2}
设 $X$ 是域 $k$ 上一个\Term{dimension}{维数} $n\geq1$ 的\Term{smooth}{光滑}\Term{complete-intersection}{完全交}。固定一个整数
\[
 d\geq1+n.
 \tag{4.2.1}
\]
记
\[
 X.H_U\longrightarrow U
 \tag{4.2.2}
\]
为 $X$ 与 $d$ 次\Term{hypersurface}{超曲面}所得、\Term{dimension}{维数}为 $n-1$ 且\Term{smooth}{光滑}的\Term{section}{截面}的\Term{universal-family}{泛族}。存在一个\Term{dense-open}{稠密开集} $T\subset U$，使得
\[
 \text{对 }t\in T,\qquad
 H^{n-1}(X.H_{\overline t},\mathbf Q_\ell)
 \ne N^1H^{n-1}(X.H_{\overline t},\mathbf Q_\ell).
 \tag{4.2.3}
\]
\end{SGAresult}

\textbf{参考文献}
\begin{enumerate}
 \item[{[1]}] J.-P. Serre，\emph{Faisceaux algébriques cohérents}（\Term{coherent-algebraic-sheaf}{凝聚代数层}），Annals of Math.，61，1955，第 197--278 页。
 \item[{[2]}] J.-P. Serre，\emph{Sur la topologie des variétés algébriques en caractéristique p}（论\Term{characteristic}{特征} $p$ 的\Term{algebraic-variety}{代数簇}的\Term{topology}{拓扑}），载 \emph{Symposium International de topologia algebraica}，Mexico，1958。
\end{enumerate}
